% ============================================================
%  Niels Treschow, "Almindelig Logik" (Kiøbenhavn: Fr. Brummer, 1813)
%  TRANSLATION (English) — Hans Halvorson
%  SELECTION: §§ 26–34, printed pp. 91–105.
%
%  Source of truth: transcription.tex. Mirrors it 1:1 (same 15 printed-page
%  markers, same 9 section headings, same emphasis spans).
%
%  TERMINOLOGY — continuous with ../enslige-ting/translation.tex:
%    enslig / enslige Ting / det Enslige -> singular / singular things
%    Individ / Individualitet            -> individual / individuality
%    besynderlig                         -> particular (18c sense)
%    Slægt / Art / Afart                 -> genus / species / variety
%    Begreb / Kiendemærke                -> concept / criterion
%    Beskrivelse                         -> description
%    Omfang (Extension)                  -> extension
%    Indhold / Stof (Comprehension)      -> comprehension
%    Deling                              -> division (generic)
%    Udstykning (partitio)               -> partition
%    Inddeling (divisio)                 -> division (logical)
%    Adskillelse (distinctio)            -> distinction
%    Heelt / Dele                        -> whole / parts
%    real / ideal                        -> real / ideal
%    Grundsetning                        -> fundamental principle
%  Treschow's Latin tags (partitio, divisio, distinctio) are kept as printed.
% ============================================================
\documentclass[12pt,a4paper]{article}
\usepackage[utf8]{inputenc}
\usepackage[T1]{fontenc}
\usepackage{libertinus}
\usepackage{libertinust1math}
\usepackage[english]{babel}
\usepackage[protrusion=true,expansion=true]{microtype}
\usepackage{setspace}
\usepackage[margin=1.2in]{geometry}
\usepackage{fancyhdr}
\usepackage{hyperref}

\hypersetup{
  pdftitle={General Logic (1813) — §§ 26–34},
  pdfauthor={Niels Treschow},
  hidelinks
}

\pagestyle{fancy}
\fancyhf{}
\rhead{Treschow, \emph{General Logic} (1813), §§ 26--34}
\cfoot{\thepage}

\setstretch{1.3}

\newcommand{\logsec}[2]{%
  \bigskip\begin{center}{\large #1}\\[4pt]{\emph{#2}}\end{center}\smallskip}

\title{\textbf{General Logic}\\[6pt]
  \large §§ 26--34 (printed pp.\ 91--105)\\[4pt]
  \normalsize The close of the treatment of definitions, and\\
  \large Fourth Division: On Divisions}
\author{Niels Treschow\\[4pt]
  \small Translated from the Danish by Hans Halvorson}
\date{\small Copenhagen: Fr.\ Brummer, 1813}

\begin{document}
\maketitle
\thispagestyle{fancy}

\bigskip

% === TRANSLATION BEGINS (printed p. 91), mid-sentence in § 25 ===

% printed p. 91
\noindent
[\dots] and sometimes also its necessity---as that by substance one must
understand the ultimate inward and real ground of all the constitutive
properties in an object: \emph{and partly that the criteria adduced occur not
only here, but likewise in all the cases designated by a certain word. The
latter is proved in turn either by induction from several examples or by
deduction,} when from a generally admitted criterion one derives that which, in
a real definition, one would rather lay down as the ground. For example, from
the fact that virtue is by ordinary usage ascribed neither to God, nor to
angels, nor to the saints, it may be concluded that virtue is moral perfection
in struggle with sensibility and a refractory nature.

\logsec{26.}{Word-explanations.}

\noindent
\emph{Word-explanations}---not the same as nominal definitions, § 22---are those
which elucidate a word's signification either by other synonymous words or by
the word's composition and origin: for example, those of philosophy or of
ontology. \emph{The aim of this is partly to prove the agreement of the
latter with the word's original and customary meaning.}
% sic: only one "partly" in the Danish, though the following "The latter"
% presupposes a second limb
The latter is especially feasible with derived and compound words. For since
those who first formed them presumably conformed therein to the natural
kinship
% printed p. 92
they thought they discovered between the signification of the root-words and the
new concept they wished to express, one may readily infer from the former what
they had in mind with the compound and derived ones. Etymology accordingly
serves to find out what words properly signified in the older language, and what
they may signify according to their origin. Root-words, on the other hand, one
cannot explain otherwise than by showing, through comparison with the derived
ones, what the former in all likelihood signified, in case this is now doubtful
or forgotten. For that either the whole word, or each letter and syllable,
should have its own meaning is scarcely possible to prove. Besides, this sort of
explanation serves to show the human understanding's own progress in the
formation of certain concepts---for example, of fate, of God, and of spiritual
beings---which is sometimes very instructive. Of less importance, to be sure, is
the etymology of such words as designate no scientific objects, but only such as
occur in daily and civil life: temples, implements, officers of authority, and
so on. Nevertheless these serve both to explain many passages in the writers,
and to acquaint us more closely with the spirit and customs of antiquity, which
can be worth the trouble of knowing.

% printed p. 93
\logsec{27.}{Which words one ought to explain.}

\noindent
Where definitions or word-explanations ought to be placed, everyone must judge
for himself according to the purpose of his exposition. \emph{In general all
principal concepts}---that is, those on which doctrines and rules are
subsequently grounded---\emph{ought to be explained exactly}, since the least
negligence here often leads to incorrect inferences. For every definition
contains a proposition that is laid down as the ground of all the considerations
one can raise about the object. It is therefore necessary that this foundation
be well laid, if all that one builds upon it is not to fall. Indistinct
concepts---though they may for the rest be both clear and correct, as those of
virtue, duty and right---do not put us in a position to answer every question
whose resolution social life and each man's own peace of mind may demand; even
though those who study jurisprudence are at first apt to regard definitions of
such well-known things as superfluous, and almost everyone imagines he knows in
advance what he ought to understand by these words. It is the same with time,
space, and other things about whose true essence the philosophers are in doubt,
while nothing seems easier for the ignorant to explain.

\emph{On the other hand, this diligence can also be overdone,} when to no
purpose one
% printed p. 94
fills textbooks with definitions of settled concepts or superfluous technical
terms, while treating the most important ones cursorily---indeed, sometimes
passing them by in complete silence. No science has been more burdened with this
kind of superfluity than metaphysics or logic itself, whence it has been carried
into other sciences as well. Good taste has sought to purge them of it. Only,
what a pity that along with such verbiage it has sometimes swept out much that
was good and solid, which it failed to notice amid the rubbish.

\logsec{28.}{Descriptions.}

\noindent
\emph{Descriptions} are properly intended for explanations of singular objects,
whose essential and abiding criteria---as with singular things generally---either
are lacking or are entirely unknown. One can therefore adduce none but
inessential and changeable ones. No inorganic body seems to have true
individuality, since no part of it is abiding, and the form itself can vary in
countless ways without one's being able to fix any point at which it ceases to
be the same---for example, Noah's Ark or Theseus's ship. In towns and buildings
all the criteria, down to the situation itself, are impermanent, as in the case
of Tyre. Human beings and the more perfect animals, at least, are beyond doubt
real individuals, which through all life's vicissitudes retain, outwardly and
% printed p. 95
inwardly, a certain indelible character; but who is in a position to determine
it? We are therefore compelled merely to describe them by name, station, age,
deeds, descent, and so forth---although all this may perish while the singular
human being or animal still remains the same. One is indeed accustomed,
especially where human beings are concerned, to add some traits of character;
but these must always remain partly uncertain, partly indeterminate, the
essential ones being easier to give occasion to divine, by indistinct hints,
than to explain exactly.

In describing general concepts, or genera and species, the investigators of
nature adduce, besides the essential or artificial criterion, also all sorts of
changeable and accidental ones. For since accidental constitutive properties
have, in part, their ground in something abiding and essential, the former may
even serve toward a more thorough knowledge of the very essence of things.
Besides, they strike the eye soonest, and are therefore easier to find out.

\bigskip
\begin{center}
  {\Large Fourth Division.}\\[6pt]
  {\large \emph{On Divisions.}}
\end{center}

\logsec{29.}{Higher and lower concepts.}

\noindent
A higher concept is a general one under which others, called lower, can be
thought as
% printed p. 96
comprehended---as human beings under animal, triangles under figure. A concept
that contains only such criteria as several things resemble one another in,
while itself being comprehended under a higher one, is called, in relation to
the latter, a \emph{species}, and the higher one a \emph{genus}. The same
general concept can therefore, in different relations, be regarded both as genus
and as species. The highest genus and the lowest species alone form an exception
to this; for the former comprehends all, without itself standing under any
higher concept, and consequently can never be regarded as a species, while the
latter comprehends only singular things, and can therefore not be regarded as a
genus.

\emph{To the higher concepts one ascribes greater extension}---extensio---or
range, because they comprehend more singulars; \emph{to the lower, more content
or matter}---comprehensio---since, besides the common criteria, they have many
peculiar ones.

\emph{Concepts are to one another either subordinate or co-ordinate.} If the
first is the case, as between genera and species, the former are called
subordinating, the latter properly subordinate. If, on the other hand, the
second, then they are either both comprehended in the same genus, as nearest
species---for example, marsh violets and pansies---or belong to different
genera---for example, the ones just named and night-violets, i.e.\ disparate
ones. \emph{Species can be subordinate to their genera either mediately or
immediately:} in the latter case
% printed p. 97
these are called their \emph{nearest}, in the former their \emph{remoter
genera}. Philosophy allows itself many leaps here---for example, from human
being to animal---which in natural history would be unpardonable.

The highest genus and the lowest species are either the \emph{absolutely highest
and lowest}---such as thing or being; human being, horse, dog, and so
forth---or \emph{relative} to this or that science. Thus body is in natural
science the highest genus, in metaphysics a much lower one. In anthropology,
where the different human races, nations, sexes, ages and so forth are treated,
the concept of a human being is the highest, which in general philosophy is the
very lowest.

\logsec{30.}{Proper and improper genera or species.}

\noindent
\emph{Genera as well as species are partly proper, partly improper or
analogical.} The criteria of the former are essential and abiding, those of the
latter merely accidental and impermanent. \emph{Although both in philosophy and
in other sciences one prefers to discover the first kind,} because it is most
incumbent upon us to know the essential properties of things, \emph{yet the
latter kind too are necessary}, when either
% printed p. 98
no other is available (§§ 21 and 23), or when they can be of use in other
respects, chiefly in the application of general theories. Thus for prudence in
dealing with human beings it is not enough to know human nature in general: the
different inclinations and circumstances by which people let themselves be
governed must likewise be known to us. Now it is nevertheless impossible to
study out the temperament or character of every singular person; and
consequently the learned have done us a service in dividing them, according to
temperament, sex, age, descent and so forth, into several improper species,
whereby the most important differences that can obtain among singular persons
are at least brought to notice. For the same reason there are entered in natural
history not only all the species regarded as abiding, but even many
varieties---for example, of pot-herbs, garden flowers and other natural
products---which can be noteworthy both in a theoretical respect and still more
for the sake of use. Many higher genera, classes and orders are themselves
improper: whence the so-called \emph{artificial systems}, as opposed to the
natural ones, have their origin.

\logsec{31.}{A whole and its parts.}

\noindent
\emph{A whole is a unity which, together with several things called the parts,
taken all together constitutes the same:}
% printed p. 99
for example, the world with all the heavenly bodies and what these contain; the
state with all its citizens; the genus with all the species belonging to it.
These parts may now either lie within and outside one another in such a way
that, by virtue of an actual or real connection in the object itself, they must
be regarded as inseparable---as in the world, and especially in organic bodies,
in which not only does each part have some influence on the rest, but no one of
them can so much as subsist by itself---or this relation may exist only in our
representations. Thus there is not among all animals of one species any actual
community that makes the one dependent on the other. Nevertheless every species
constitutes a whole. In the first case they are \emph{real}, in the second
\emph{ideal parts}. Singular beings are parts of the species, these of the
genera, and all taken together parts of the higher genus---but only in thought.
\emph{This difference is very important.} One has often confused the real
totality, which is the world, with the ideal one, which is thing or being in
general; and both again with the highest unity in principles, or the ultimate
cause of all things, on which they depend as effects. Pantheism and the
emanation-system rest chiefly upon this misunderstanding. The highest being is
neither a general concept, under which spirit and body are comprehended as
species, nor yet any whole, like the world, of which all individuals are parts.
But in order to root out such representations it is necessary
% printed p. 100
to develop these concepts exactly, though in themselves they seem clear enough.

\logsec{32.}{Logical and other divisions.}

\noindent
From this it follows that a whole can also be divided in two ways: namely into
real parts, which actually cohere or are inwardly connected---a resolution into
such may be called \emph{partition}---partitio---and into ideal ones, when a
genus is resolved into its species, which is properly, or \emph{logical
division}---divisio. In the first way one divides countries into their
provinces, the body into head, trunk and limbs, civil society into estates, and
so on. In the second, natural history and other sciences produce their classes,
orders, genera, species and varieties; of which many occur in logic too---for
example, in the division of concepts.

\emph{Distinction}---distinctio---is, by contrast, only a remark upon the
difference of several closely related concepts, which do not even always belong
to the same genus or species: for example, synonyms, such as luxury and
extravagance; sensible, shrewd, rational, and the like. How much this exactness
conduces to distinctness in concepts is clear from the foregoing.

\emph{The members or limbs of a division may be two, three or more}---dichotomy,
trichotomy, and so on.
% printed p. 101
---\emph{How many there ought to be is partly already determined by experience,
through the essential difference of the natural genera and species}, as in
natural history, where sometimes scores of species are given under the same
genus; \emph{and partly it must be determined by ourselves according to certain
fundamental principles.} But among these, some may favour dichotomy, others
trichotomy. Thus the principle that every thing has its opposite, and that all
things are thereby both born and perish, leads to dichotomy; whereas the
principle that such opposition must always be united by some third, without
which no thing could be born or subsist, leads to division into three---which
the critical and most recent philosophy of nature therefore prefer. The
sacredness or higher significance of the number three is thus acknowledged here
as well.

\emph{Divisions may be made in several respects}, or with regard to several of
the whole's inward or outward criteria, \emph{and these are then called the
ground of the division.} Since concepts, for example, have both matter and form,
they may in the first respect be divided into simple and compound; in the
second, into clear and obscure; again, as regards their origin, into pure and
empirical, and so forth. In the same ways human beings may be divided according
to sex, age, aptitude, station, and so on, each of which yields a particular
ground for new species and subspecies.

\emph{The division may be continued by subdivisions}, when the species are again
resolved into new ones. Thus one divides, for example, clear concepts into
distinct and indistinct, the first
% printed p. 102
into thorough and superficial (§ 3); animals and plants first into certain
classes and families, then into genera, species and varieties.

\logsec{33.}{The utility of divisions.}

\noindent
\emph{Divisions are very useful a) for surveying and ordering the whole
collection of knowledge one has acquired.} Without divisions it would not only
be confused, but difficult to retain, to use and to recall. To this is added the
pleasure that unity amid so great a manifoldness, and the enjoyment of our whole
riches in a general survey, can afford. \emph{This is why the tabular method has
found so many admirers in all the sciences.} It consists in setting forth, in
order but briefly, all the divisions, subdivisions and partitions of the
principal concept or object, down to the most particular; the doctrines
belonging to each concept being added in their place, but without proof. This
manner is not only an almost indispensable aid to the memory---chiefly in the
empirical sciences and disciplines---but is also useful for discerning the real
as well as the ideal connection of things. Nevertheless it can be carried too
far even in these, when it is made exclusive, and the memory is thereby enriched
more than the higher power of thought is sharpened. In the stricter or proper
% printed p. 103
sciences especially, it then leads easily to superficiality and verbiage: of
which the Lullian art affords a remarkable proof.

\emph{b) For easing and guiding the work of judgment in the application of
general concepts or rules.} For this reason the division must be continued down
to the lowest species of things. In the practical sciences it is especially
necessary to descend as deep as possible, without engaging with merely singular
cases, from which an endless prolixity would arise within the science itself. In
agriculture and horticulture one must therefore know even the most important
varieties of trees, grains and other plants, which general natural history
passes over. The statesman's knowledge must not confine itself to any merely
general knowledge of forms of government and of states. Every individual of a
state is important enough to be considered for itself. For the moralist and the
popular teacher it is not enough to know the principal species of virtues and
vices---to which the diligence of the scholastic philosophers in this matter
mostly confined itself. A general discourse on such subjects commonly remains
unfruitful, unless one also knows and understands how to apply it to the many
sorts of ways in which virtue and vice are wont to vary according to persons'
outward and inward condition.

% printed p. 104
\logsec{34.}{Rules necessary throughout.}

\noindent
In order to derive this utility from divisions, the following rules are at least
to be observed. \emph{a) The parts must be enumerated completely, and the whole
which one divides must thereby be exhausted}; for otherwise our knowledge of it
becomes not only incomplete, but can even, when one nevertheless imagines the
contrary, mislead us into false inferences. As if one were to divide salts into
alkaline and acid salts, from which it would follow that common salt and sugar
belonged either to one of the parts or to another genus. Were one to say that
all virtue is either justice or magnanimity, then both wisdom and truthfulness,
with several others, would have to be excluded. Were one likewise to divide
animals into those with two feet, four feet, and none, then those with six or
more feet, such as insects, would come to stand outside this class.

\emph{b) They must be mutually opposed, and not comprehend one another}---as
reciprocity does in causality, although Kant, under the category of relation,
makes two species of these, just as he makes infinite propositions a species of
their own, distinct from the affirmative and the negative; of which more
hereafter. For the members of the division then coincide, and no real division
takes place.

\emph{c) No leap must be made from higher to lower species; rather the nearest
must be set
% printed p. 105
beside one another:} for example, if instead of dividing figures into
rectilinear and curvilinear, or regular and irregular, one were at once to
divide them into circles, triangles, squares and polygons, one would offend at
one stroke both against this rule and against the first. No better is the common
division of forms of government into monarchic, democratic, aristocratic and
despotic---which even combines the aforesaid faults with the following one.

\emph{d) No foreign species ought to be introduced into the division.} A rule
whose correctness is indeed self-evident, but which is nevertheless often
transgressed for want of a distinct concept of the whole---as when either
despotism, or what is still worse, ochlocracy, indeed even anarchy, is set
beside the regular forms of government.

\logsec{35.}{Other rules.}

\noindent
Not equally necessary, but nevertheless important where they can and ought to be
observed, are the following. \emph{a) The ground of the division ought to be
taken from the essential and abiding criteria of things.} If none such are to be
found, one must choose, among several others, those that come nearest to them.
This rule is actually followed by all the more recent describers of nature and
anthropologists, although there may again be doubt as to which
% printed p. 106
constitutive properties are the most abiding and the most closely akin to the
essence itself. The older writers, by contrast, divided plants, for example,
according to size---which is nevertheless highly variable---into trees, shrubs
and herbs; and human beings according to the countries they inhabited or the
language they spoke, into Greeks and barbarians. It would be just as incorrect
to divide states into empires, kingdoms, duchies or principalities, and so
forth. Every genus can be divided either by inward or by outward criteria. Among
the former, \emph{form, which properly determines the essence, is more important
than matter; among the latter, cause and origin are more important than effect
and object.} Only so far as they all have some influence upon the form of things
do they come more or less into consideration. Where one must, in order to avoid
excessive prolixity, confine oneself to a single classification---as of the
duties in moral doctrine, or of the sciences in an encyclopaedia---this rule is
especially necessary to observe. Yet even here one is often obliged to conform
to ordinary usage, which presupposes a certain classification as already made,
and which one sometimes cannot alter without introducing a wholly new
nomenclature---something that seldom succeeds.

\emph{b) The ground of division must be indicated, chiefly when there are
several;} since otherwise one can expect no determinate answer, nor discern how
the various species are co-ordinated with one another, or in what mutual
relation they stand. Never%
% printed p. 107
theless the division may be correct without this being done. Sometimes, too,
when there is only one principal division, it is superfluous.

\emph{c) As regards the number of members, one ought indeed to follow nature;
but where nature leaves us free to choose}---as when the objects of concepts are
given not by experience but by the higher power of thought
itself---\emph{one must prefer dichotomy, and rather make several subdivisions
out of it.} Thus forms of government, for example, may be divided into regular
and irregular. In the first, either one or several govern; and where several,
either all or some.

The ground of this Ramist rule is clear: namely, that otherwise one can have no
certainty that the division has due completeness. In every other kind of
division the members may indeed be opposed, but not directly contradictory, and
only in the latter case do they exclude the possibility of some further member.
As for the ground for trichotomy adduced earlier in § 32, it must be remarked
that conflicting things can on the whole be united, but not if they are
contradictory. One will therefore also find that trichotomy in this case always
fails; while in the other it does not exhaust the whole concept, because of
conflicting things a far greater number may be given---as will be shown more
distinctly in what follows.
% printed p. 108
The three dimensions of space prove nothing against this. For properly they
are, like the directions of motion, innumerable. But according to the position
of our body---not according to any universally valid principle---we refer them
all to three principal kinds: namely fore and aft, right and left, up and down;
of which each has its opposite.

How much arbitrariness both trichotomy and every other mode of classification
gives occasion to is shown by the whole history of philosophy, especially the
doctrine of the categories and the most recent philosophy of nature. Perhaps one
will find here the reason why the opinion about the excellence or sacredness of
certain numbers has found so many adherents among the philosophers. For when one
has no certain rule to follow, or has once left the beaten path, only
bewildering byways remain to guide our step.

Should the whole be so composed that among the conflicting predicates the one
belongs to the one part and the other to the other, then dichotomy may still be
retained, in that one divides the species---for example, forms of government, or
feelings---into pure and mixed, the latter being then again easily distinguished
into two.

% [End of the First Main Part; p. 109 begins "Andet Hovedstykke. Om Sandhed og
%  Vished" (Second Main Part, On Truth and Certainty), 1st Division, § 36
%  "Hvad er Sandhed?" — not included.]

\begin{center}\rule{0.25\linewidth}{0.4pt}\end{center}

\end{document}
